% Nomenclature — Passage du 3D au 1D
% Préambule : \usepackage{nomencl} et \makenomenclature
% Document  : % Nomenclature — Passage du 3D au 1D
% Préambule : \usepackage{nomencl} et \makenomenclature
% Document  : % Nomenclature — Passage du 3D au 1D
% Préambule : \usepackage{nomencl} et \makenomenclature
% Document  : % Nomenclature — Passage du 3D au 1D
% Préambule : \usepackage{nomencl} et \makenomenclature
% Document  : \input{nomenclature} puis \printnomenclature
% Compiler  : makeindex main.nlo -s nomencl.ist -o main.nls

% ------------------------------------------------------------
% LETTRES LATINES
% ------------------------------------------------------------
\nomenclature[Aav]{$a_v$}{Aire interfaciale spécifique volumique $= A_{sf}/V$ ($\text{m}^{-1}$)}
\nomenclature[AAsf]{$A_{sf}$}{Aire de l'interface fluide-solide ($\text{m}^2$)}
\nomenclature[Ab]{$\mathbf{b}$}{Vecteur de fermeture pour la concentration (m)}
\nomenclature[ACi]{$C_i$}{Concentration de l'espèce $i$ ($\text{mol m}^{-3}$)}
\nomenclature[ACref]{$C_\text{ref}$}{Concentration de référence ($\text{mol m}^{-3}$)}
\nomenclature[Ad]{$\mathbf{d}$}{Vecteur de fermeture pour le potentiel électrique (m)}
\nomenclature[ADa]{$Da$}{Nombre de Damköhler $= j\,n\,l_e / (F D_i C_\text{ref})$}
\nomenclature[ADi]{$D_i$}{Coefficient de diffusion de l'espèce $i$ ($\text{m}^2\text{s}^{-1}$)}
\nomenclature[ADei]{$D^{eff}_i$}{Coefficient de diffusion effectif de l'espèce $i$ ($\text{m}^2\text{s}^{-1}$)}
\nomenclature[ADmig]{$D^{eff,\text{mig}}_i$}{Coefficient de migration effectif de l'espèce $i$ ($\text{m}^2\text{s}^{-1}$)}
\nomenclature[ADdisp]{$D_\text{disp}$}{Coefficient de dispersion ($\text{m}^2\text{s}^{-1}$)}
\nomenclature[ADtot]{$D^\text{tot}_i$}{Coefficient de diffusion total, diffusion et dispersion ($\text{m}^2\text{s}^{-1}$)}
\nomenclature[Ae]{$\mathbf{e}$}{Vecteur de fermeture pour le potentiel solide (m)}
\nomenclature[AEeq]{$E_\text{eq}$}{Potentiel d'équilibre (V)}
\nomenclature[AE]{$E$}{Potentiel imposé au collecteur de courant (V)}
\nomenclature[AF]{$F$}{Constante de Faraday $= 96\,485$ C mol$^{-1}$}
\nomenclature[Ag]{$g$}{Grandeur physique quelconque définie dans le fluide}
\nomenclature[Aj]{$j$}{Densité de courant interfacial ($\text{A m}^{-2}$)}
\nomenclature[Aj0]{$j_0$}{Densité de courant d'échange ($\text{A m}^{-2}$)}
\nomenclature[Ak]{$k$}{Constante de réaction du premier ordre ($\text{m s}^{-1}$)}
\nomenclature[Akv]{$k_v$}{Constante de réaction volumique $= \nu_i a_v k(\eta)/(n_e F)$ ($\text{s}^{-1}$)}
%\nomenclature[AL]{$L$}{Longueur caractéristique macroscopique (m)}
\nomenclature[Ale]{$l_e$}{Épaisseur de l'électrode (m)}
\nomenclature[Alp]{$\ell_p$}{Longueur caractéristique à l'échelle du pore (m)}
\nomenclature[Ana]{$\mathbf{n}$}{Vecteur normal unitaire (pointant du fluide vers le solide)}
\nomenclature[Anb]{$\mathbf{n}_s$}{Vecteur normal unitaire sortant de la phase solide ; $\mathbf{n}_s = -\mathbf{n}$}
\nomenclature[Ane]{$n_e$}{Nombre d'électrons échangés lors de la réaction}
\nomenclature[ANi]{$N_i$}{Flux molaire de l'espèce $i$ ($\text{mol m}^{-2}\text{s}^{-1}$)}
\nomenclature[APe]{$Pe$}{Nombre de Péclet $= v\,l_e/D_\text{ref}$}
\nomenclature[AR]{$R$}{Constante des gaz parfaits $= 8.314$ J mol$^{-1}$ K$^{-1}$}
\nomenclature[ARi]{$R_i$}{Taux de réaction interfaciale de l'espèce $i$ ($\text{mol m}^{-2}\text{s}^{-1}$)}
\nomenclature[AT]{$T$}{Température (K)}
\nomenclature[Av]{$v$}{Vitesse du fluide ($\text{m s}^{-1}$)}
\nomenclature[AV]{$V$}{Volume total de contrôle, fluide et solide ($\text{m}^3$)}
\nomenclature[AVl]{$V_l$}{Volume occupé par le fluide ($\text{m}^3$)}
\nomenclature[AVs]{$V_s$}{Volume occupé par le solide ($\text{m}^3$)}
\nomenclature[Azi]{$z_i$}{Valence de l'espèce $i$}

% ------------------------------------------------------------
% LETTRES GRECQUES
% ------------------------------------------------------------
\nomenclature[Baa]{$\alpha_a,\, \alpha_c$}{Coefficients de transfert anodique et cathodique}
\nomenclature[Bdelta]{$\bar{\delta}$}{Longueur de Debye adimensionnée $= \lambda_0/\ell_p$}
\nomenclature[Bepsl]{$\varepsilon_l$}{Porosité du milieu $= V_l/V$}
\nomenclature[Bepss]{$\varepsilon_s$}{Fraction volumique du solide $= V_s/V$}
\nomenclature[Bepse]{$\varepsilon^e$}{Permittivité effective du milieu poreux}
\nomenclature[Bepsr]{$\varepsilon_r$}{Permittivité relative du fluide}
\nomenclature[Beps0]{$\varepsilon_0$}{Permittivité du vide $= 8.854\times10^{-12}$ F m$^{-1}$}
\nomenclature[BetaE]{$\eta_E$}{Facteur d'efficacité}
\nomenclature[Beta]{$\eta$}{Surtension (V)}
\nomenclature[Blambda]{$\lambda_0$}{Longueur de Debye (m)}
\nomenclature[Bnu]{$\nu_i$}{Coefficient stœchiométrique de l'espèce $i$}
\nomenclature[Bphi]{$\varphi$}{Module de Thiele microscopique}
%\nomenclature[Bphimacro]{$\varphi_\text{macro}$}{Module de Thiele macroscopique}
\nomenclature[Bphil]{$\varphi_l$}{Potentiel électrique dans le fluide (V)}
\nomenclature[Bphis]{$\varphi_s$}{Potentiel électrique dans le solide (V)}
\nomenclature[Bsigmas]{$\sigma_s$}{Conductivité électrique du solide ($\text{S m}^{-1}$)}
\nomenclature[Bsigmaq]{$\sigma_q$}{Densité de charge surfacique à l'interface ($\text{C m}^{-2}$)}
\nomenclature[Bsigmae]{$\sigma^e$}{Conductivité effective du milieu poreux ($\text{S m}^{-1}$)}

% ------------------------------------------------------------
% OPÉRATEURS ET MOYENNES
% ------------------------------------------------------------
\nomenclature[Cavg]{$\langle g \rangle$}{Moyenne superficielle de $g$}
\nomenclature[Cavgl]{$\langle g \rangle^l$}{Moyenne intrinsèque de $g$ dans le fluide}
\nomenclature[Ctilde]{$\tilde{g}$}{Fluctuation de $g$ à l'échelle du pore} puis \printnomenclature
% Compiler  : makeindex main.nlo -s nomencl.ist -o main.nls

% ------------------------------------------------------------
% LETTRES LATINES
% ------------------------------------------------------------
\nomenclature[Aav]{$a_v$}{Aire interfaciale spécifique volumique $= A_{sf}/V$ ($\text{m}^{-1}$)}
\nomenclature[AAsf]{$A_{sf}$}{Aire de l'interface fluide-solide ($\text{m}^2$)}
\nomenclature[Ab]{$\mathbf{b}$}{Vecteur de fermeture pour la concentration (m)}
\nomenclature[ACi]{$C_i$}{Concentration de l'espèce $i$ ($\text{mol m}^{-3}$)}
\nomenclature[ACref]{$C_\text{ref}$}{Concentration de référence ($\text{mol m}^{-3}$)}
\nomenclature[Ad]{$\mathbf{d}$}{Vecteur de fermeture pour le potentiel électrique (m)}
\nomenclature[ADa]{$Da$}{Nombre de Damköhler $= j\,n\,l_e / (F D_i C_\text{ref})$}
\nomenclature[ADi]{$D_i$}{Coefficient de diffusion de l'espèce $i$ ($\text{m}^2\text{s}^{-1}$)}
\nomenclature[ADei]{$D^{eff}_i$}{Coefficient de diffusion effectif de l'espèce $i$ ($\text{m}^2\text{s}^{-1}$)}
\nomenclature[ADmig]{$D^{eff,\text{mig}}_i$}{Coefficient de migration effectif de l'espèce $i$ ($\text{m}^2\text{s}^{-1}$)}
\nomenclature[ADdisp]{$D_\text{disp}$}{Coefficient de dispersion ($\text{m}^2\text{s}^{-1}$)}
\nomenclature[ADtot]{$D^\text{tot}_i$}{Coefficient de diffusion total, diffusion et dispersion ($\text{m}^2\text{s}^{-1}$)}
\nomenclature[Ae]{$\mathbf{e}$}{Vecteur de fermeture pour le potentiel solide (m)}
\nomenclature[AEeq]{$E_\text{eq}$}{Potentiel d'équilibre (V)}
\nomenclature[AE]{$E$}{Potentiel imposé au collecteur de courant (V)}
\nomenclature[AF]{$F$}{Constante de Faraday $= 96\,485$ C mol$^{-1}$}
\nomenclature[Ag]{$g$}{Grandeur physique quelconque définie dans le fluide}
\nomenclature[Aj]{$j$}{Densité de courant interfacial ($\text{A m}^{-2}$)}
\nomenclature[Aj0]{$j_0$}{Densité de courant d'échange ($\text{A m}^{-2}$)}
\nomenclature[Ak]{$k$}{Constante de réaction du premier ordre ($\text{m s}^{-1}$)}
\nomenclature[Akv]{$k_v$}{Constante de réaction volumique $= \nu_i a_v k(\eta)/(n_e F)$ ($\text{s}^{-1}$)}
%\nomenclature[AL]{$L$}{Longueur caractéristique macroscopique (m)}
\nomenclature[Ale]{$l_e$}{Épaisseur de l'électrode (m)}
\nomenclature[Alp]{$\ell_p$}{Longueur caractéristique à l'échelle du pore (m)}
\nomenclature[Ana]{$\mathbf{n}$}{Vecteur normal unitaire (pointant du fluide vers le solide)}
\nomenclature[Anb]{$\mathbf{n}_s$}{Vecteur normal unitaire sortant de la phase solide ; $\mathbf{n}_s = -\mathbf{n}$}
\nomenclature[Ane]{$n_e$}{Nombre d'électrons échangés lors de la réaction}
\nomenclature[ANi]{$N_i$}{Flux molaire de l'espèce $i$ ($\text{mol m}^{-2}\text{s}^{-1}$)}
\nomenclature[APe]{$Pe$}{Nombre de Péclet $= v\,l_e/D_\text{ref}$}
\nomenclature[AR]{$R$}{Constante des gaz parfaits $= 8.314$ J mol$^{-1}$ K$^{-1}$}
\nomenclature[ARi]{$R_i$}{Taux de réaction interfaciale de l'espèce $i$ ($\text{mol m}^{-2}\text{s}^{-1}$)}
\nomenclature[AT]{$T$}{Température (K)}
\nomenclature[Av]{$v$}{Vitesse du fluide ($\text{m s}^{-1}$)}
\nomenclature[AV]{$V$}{Volume total de contrôle, fluide et solide ($\text{m}^3$)}
\nomenclature[AVl]{$V_l$}{Volume occupé par le fluide ($\text{m}^3$)}
\nomenclature[AVs]{$V_s$}{Volume occupé par le solide ($\text{m}^3$)}
\nomenclature[Azi]{$z_i$}{Valence de l'espèce $i$}

% ------------------------------------------------------------
% LETTRES GRECQUES
% ------------------------------------------------------------
\nomenclature[Baa]{$\alpha_a,\, \alpha_c$}{Coefficients de transfert anodique et cathodique}
\nomenclature[Bdelta]{$\bar{\delta}$}{Longueur de Debye adimensionnée $= \lambda_0/\ell_p$}
\nomenclature[Bepsl]{$\varepsilon_l$}{Porosité du milieu $= V_l/V$}
\nomenclature[Bepss]{$\varepsilon_s$}{Fraction volumique du solide $= V_s/V$}
\nomenclature[Bepse]{$\varepsilon^e$}{Permittivité effective du milieu poreux}
\nomenclature[Bepsr]{$\varepsilon_r$}{Permittivité relative du fluide}
\nomenclature[Beps0]{$\varepsilon_0$}{Permittivité du vide $= 8.854\times10^{-12}$ F m$^{-1}$}
\nomenclature[BetaE]{$\eta_E$}{Facteur d'efficacité}
\nomenclature[Beta]{$\eta$}{Surtension (V)}
\nomenclature[Blambda]{$\lambda_0$}{Longueur de Debye (m)}
\nomenclature[Bnu]{$\nu_i$}{Coefficient stœchiométrique de l'espèce $i$}
\nomenclature[Bphi]{$\varphi$}{Module de Thiele microscopique}
%\nomenclature[Bphimacro]{$\varphi_\text{macro}$}{Module de Thiele macroscopique}
\nomenclature[Bphil]{$\varphi_l$}{Potentiel électrique dans le fluide (V)}
\nomenclature[Bphis]{$\varphi_s$}{Potentiel électrique dans le solide (V)}
\nomenclature[Bsigmas]{$\sigma_s$}{Conductivité électrique du solide ($\text{S m}^{-1}$)}
\nomenclature[Bsigmaq]{$\sigma_q$}{Densité de charge surfacique à l'interface ($\text{C m}^{-2}$)}
\nomenclature[Bsigmae]{$\sigma^e$}{Conductivité effective du milieu poreux ($\text{S m}^{-1}$)}

% ------------------------------------------------------------
% OPÉRATEURS ET MOYENNES
% ------------------------------------------------------------
\nomenclature[Cavg]{$\langle g \rangle$}{Moyenne superficielle de $g$}
\nomenclature[Cavgl]{$\langle g \rangle^l$}{Moyenne intrinsèque de $g$ dans le fluide}
\nomenclature[Ctilde]{$\tilde{g}$}{Fluctuation de $g$ à l'échelle du pore} puis \printnomenclature
% Compiler  : makeindex main.nlo -s nomencl.ist -o main.nls

% ------------------------------------------------------------
% LETTRES LATINES
% ------------------------------------------------------------
\nomenclature[Aav]{$a_v$}{Aire interfaciale spécifique volumique $= A_{sf}/V$ ($\text{m}^{-1}$)}
\nomenclature[AAsf]{$A_{sf}$}{Aire de l'interface fluide-solide ($\text{m}^2$)}
\nomenclature[Ab]{$\mathbf{b}$}{Vecteur de fermeture pour la concentration (m)}
\nomenclature[ACi]{$C_i$}{Concentration de l'espèce $i$ ($\text{mol m}^{-3}$)}
\nomenclature[ACref]{$C_\text{ref}$}{Concentration de référence ($\text{mol m}^{-3}$)}
\nomenclature[Ad]{$\mathbf{d}$}{Vecteur de fermeture pour le potentiel électrique (m)}
\nomenclature[ADa]{$Da$}{Nombre de Damköhler $= j\,n\,l_e / (F D_i C_\text{ref})$}
\nomenclature[ADi]{$D_i$}{Coefficient de diffusion de l'espèce $i$ ($\text{m}^2\text{s}^{-1}$)}
\nomenclature[ADei]{$D^{eff}_i$}{Coefficient de diffusion effectif de l'espèce $i$ ($\text{m}^2\text{s}^{-1}$)}
\nomenclature[ADmig]{$D^{eff,\text{mig}}_i$}{Coefficient de migration effectif de l'espèce $i$ ($\text{m}^2\text{s}^{-1}$)}
\nomenclature[ADdisp]{$D_\text{disp}$}{Coefficient de dispersion ($\text{m}^2\text{s}^{-1}$)}
\nomenclature[ADtot]{$D^\text{tot}_i$}{Coefficient de diffusion total, diffusion et dispersion ($\text{m}^2\text{s}^{-1}$)}
\nomenclature[Ae]{$\mathbf{e}$}{Vecteur de fermeture pour le potentiel solide (m)}
\nomenclature[AEeq]{$E_\text{eq}$}{Potentiel d'équilibre (V)}
\nomenclature[AE]{$E$}{Potentiel imposé au collecteur de courant (V)}
\nomenclature[AF]{$F$}{Constante de Faraday $= 96\,485$ C mol$^{-1}$}
\nomenclature[Ag]{$g$}{Grandeur physique quelconque définie dans le fluide}
\nomenclature[Aj]{$j$}{Densité de courant interfacial ($\text{A m}^{-2}$)}
\nomenclature[Aj0]{$j_0$}{Densité de courant d'échange ($\text{A m}^{-2}$)}
\nomenclature[Ak]{$k$}{Constante de réaction du premier ordre ($\text{m s}^{-1}$)}
\nomenclature[Akv]{$k_v$}{Constante de réaction volumique $= \nu_i a_v k(\eta)/(n_e F)$ ($\text{s}^{-1}$)}
%\nomenclature[AL]{$L$}{Longueur caractéristique macroscopique (m)}
\nomenclature[Ale]{$l_e$}{Épaisseur de l'électrode (m)}
\nomenclature[Alp]{$\ell_p$}{Longueur caractéristique à l'échelle du pore (m)}
\nomenclature[Ana]{$\mathbf{n}$}{Vecteur normal unitaire (pointant du fluide vers le solide)}
\nomenclature[Anb]{$\mathbf{n}_s$}{Vecteur normal unitaire sortant de la phase solide ; $\mathbf{n}_s = -\mathbf{n}$}
\nomenclature[Ane]{$n_e$}{Nombre d'électrons échangés lors de la réaction}
\nomenclature[ANi]{$N_i$}{Flux molaire de l'espèce $i$ ($\text{mol m}^{-2}\text{s}^{-1}$)}
\nomenclature[APe]{$Pe$}{Nombre de Péclet $= v\,l_e/D_\text{ref}$}
\nomenclature[AR]{$R$}{Constante des gaz parfaits $= 8.314$ J mol$^{-1}$ K$^{-1}$}
\nomenclature[ARi]{$R_i$}{Taux de réaction interfaciale de l'espèce $i$ ($\text{mol m}^{-2}\text{s}^{-1}$)}
\nomenclature[AT]{$T$}{Température (K)}
\nomenclature[Av]{$v$}{Vitesse du fluide ($\text{m s}^{-1}$)}
\nomenclature[AV]{$V$}{Volume total de contrôle, fluide et solide ($\text{m}^3$)}
\nomenclature[AVl]{$V_l$}{Volume occupé par le fluide ($\text{m}^3$)}
\nomenclature[AVs]{$V_s$}{Volume occupé par le solide ($\text{m}^3$)}
\nomenclature[Azi]{$z_i$}{Valence de l'espèce $i$}

% ------------------------------------------------------------
% LETTRES GRECQUES
% ------------------------------------------------------------
\nomenclature[Baa]{$\alpha_a,\, \alpha_c$}{Coefficients de transfert anodique et cathodique}
\nomenclature[Bdelta]{$\bar{\delta}$}{Longueur de Debye adimensionnée $= \lambda_0/\ell_p$}
\nomenclature[Bepsl]{$\varepsilon_l$}{Porosité du milieu $= V_l/V$}
\nomenclature[Bepss]{$\varepsilon_s$}{Fraction volumique du solide $= V_s/V$}
\nomenclature[Bepse]{$\varepsilon^e$}{Permittivité effective du milieu poreux}
\nomenclature[Bepsr]{$\varepsilon_r$}{Permittivité relative du fluide}
\nomenclature[Beps0]{$\varepsilon_0$}{Permittivité du vide $= 8.854\times10^{-12}$ F m$^{-1}$}
\nomenclature[BetaE]{$\eta_E$}{Facteur d'efficacité}
\nomenclature[Beta]{$\eta$}{Surtension (V)}
\nomenclature[Blambda]{$\lambda_0$}{Longueur de Debye (m)}
\nomenclature[Bnu]{$\nu_i$}{Coefficient stœchiométrique de l'espèce $i$}
\nomenclature[Bphi]{$\varphi$}{Module de Thiele microscopique}
%\nomenclature[Bphimacro]{$\varphi_\text{macro}$}{Module de Thiele macroscopique}
\nomenclature[Bphil]{$\varphi_l$}{Potentiel électrique dans le fluide (V)}
\nomenclature[Bphis]{$\varphi_s$}{Potentiel électrique dans le solide (V)}
\nomenclature[Bsigmas]{$\sigma_s$}{Conductivité électrique du solide ($\text{S m}^{-1}$)}
\nomenclature[Bsigmaq]{$\sigma_q$}{Densité de charge surfacique à l'interface ($\text{C m}^{-2}$)}
\nomenclature[Bsigmae]{$\sigma^e$}{Conductivité effective du milieu poreux ($\text{S m}^{-1}$)}

% ------------------------------------------------------------
% OPÉRATEURS ET MOYENNES
% ------------------------------------------------------------
\nomenclature[Cavg]{$\langle g \rangle$}{Moyenne superficielle de $g$}
\nomenclature[Cavgl]{$\langle g \rangle^l$}{Moyenne intrinsèque de $g$ dans le fluide}
\nomenclature[Ctilde]{$\tilde{g}$}{Fluctuation de $g$ à l'échelle du pore} puis \printnomenclature
% Compiler  : makeindex main.nlo -s nomencl.ist -o main.nls

% ------------------------------------------------------------
% LETTRES LATINES
% ------------------------------------------------------------
\nomenclature[Aav]{$a_v$}{Aire interfaciale spécifique volumique $= A_{sf}/V$ ($\text{m}^{-1}$)}
\nomenclature[AAsf]{$A_{sf}$}{Aire de l'interface fluide-solide ($\text{m}^2$)}
\nomenclature[Ab]{$\mathbf{b}$}{Vecteur de fermeture pour la concentration (m)}
\nomenclature[ACi]{$C_i$}{Concentration de l'espèce $i$ ($\text{mol m}^{-3}$)}
\nomenclature[ACref]{$C_\text{ref}$}{Concentration de référence ($\text{mol m}^{-3}$)}
\nomenclature[Ad]{$\mathbf{d}$}{Vecteur de fermeture pour le potentiel électrique (m)}
\nomenclature[ADa]{$Da$}{Nombre de Damköhler $= j\,n\,l_e / (F D_i C_\text{ref})$}
\nomenclature[ADi]{$D_i$}{Coefficient de diffusion de l'espèce $i$ ($\text{m}^2\text{s}^{-1}$)}
\nomenclature[ADei]{$D^{eff}_i$}{Coefficient de diffusion effectif de l'espèce $i$ ($\text{m}^2\text{s}^{-1}$)}
\nomenclature[ADmig]{$D^{eff,\text{mig}}_i$}{Coefficient de migration effectif de l'espèce $i$ ($\text{m}^2\text{s}^{-1}$)}
\nomenclature[ADdisp]{$D_\text{disp}$}{Coefficient de dispersion ($\text{m}^2\text{s}^{-1}$)}
\nomenclature[ADtot]{$D^\text{tot}_i$}{Coefficient de diffusion total, diffusion et dispersion ($\text{m}^2\text{s}^{-1}$)}
\nomenclature[Ae]{$\mathbf{e}$}{Vecteur de fermeture pour le potentiel solide (m)}
\nomenclature[AEeq]{$E_\text{eq}$}{Potentiel d'équilibre (V)}
\nomenclature[AE]{$E$}{Potentiel imposé au collecteur de courant (V)}
\nomenclature[AF]{$F$}{Constante de Faraday $= 96\,485$ C mol$^{-1}$}
\nomenclature[Ag]{$g$}{Grandeur physique quelconque définie dans le fluide}
\nomenclature[Aj]{$j$}{Densité de courant interfacial ($\text{A m}^{-2}$)}
\nomenclature[Aj0]{$j_0$}{Densité de courant d'échange ($\text{A m}^{-2}$)}
\nomenclature[Ak]{$k$}{Constante de réaction du premier ordre ($\text{m s}^{-1}$)}
\nomenclature[Akv]{$k_v$}{Constante de réaction volumique $= \nu_i a_v k(\eta)/(n_e F)$ ($\text{s}^{-1}$)}
%\nomenclature[AL]{$L$}{Longueur caractéristique macroscopique (m)}
\nomenclature[Ale]{$l_e$}{Épaisseur de l'électrode (m)}
\nomenclature[Alp]{$\ell_p$}{Longueur caractéristique à l'échelle du pore (m)}
\nomenclature[Ana]{$\mathbf{n}$}{Vecteur normal unitaire (pointant du fluide vers le solide)}
\nomenclature[Anb]{$\mathbf{n}_s$}{Vecteur normal unitaire sortant de la phase solide ; $\mathbf{n}_s = -\mathbf{n}$}
\nomenclature[Ane]{$n_e$}{Nombre d'électrons échangés lors de la réaction}
\nomenclature[ANi]{$N_i$}{Flux molaire de l'espèce $i$ ($\text{mol m}^{-2}\text{s}^{-1}$)}
\nomenclature[APe]{$Pe$}{Nombre de Péclet $= v\,l_e/D_\text{ref}$}
\nomenclature[AR]{$R$}{Constante des gaz parfaits $= 8.314$ J mol$^{-1}$ K$^{-1}$}
\nomenclature[ARi]{$R_i$}{Taux de réaction interfaciale de l'espèce $i$ ($\text{mol m}^{-2}\text{s}^{-1}$)}
\nomenclature[AT]{$T$}{Température (K)}
\nomenclature[Av]{$v$}{Vitesse du fluide ($\text{m s}^{-1}$)}
\nomenclature[AV]{$V$}{Volume total de contrôle, fluide et solide ($\text{m}^3$)}
\nomenclature[AVl]{$V_l$}{Volume occupé par le fluide ($\text{m}^3$)}
\nomenclature[AVs]{$V_s$}{Volume occupé par le solide ($\text{m}^3$)}
\nomenclature[Azi]{$z_i$}{Valence de l'espèce $i$}

% ------------------------------------------------------------
% LETTRES GRECQUES
% ------------------------------------------------------------
\nomenclature[Baa]{$\alpha_a,\, \alpha_c$}{Coefficients de transfert anodique et cathodique}
\nomenclature[Bdelta]{$\bar{\delta}$}{Longueur de Debye adimensionnée $= \lambda_0/\ell_p$}
\nomenclature[Bepsl]{$\varepsilon_l$}{Porosité du milieu $= V_l/V$}
\nomenclature[Bepss]{$\varepsilon_s$}{Fraction volumique du solide $= V_s/V$}
\nomenclature[Bepse]{$\varepsilon^e$}{Permittivité effective du milieu poreux}
\nomenclature[Bepsr]{$\varepsilon_r$}{Permittivité relative du fluide}
\nomenclature[Beps0]{$\varepsilon_0$}{Permittivité du vide $= 8.854\times10^{-12}$ F m$^{-1}$}
\nomenclature[BetaE]{$\eta_E$}{Facteur d'efficacité}
\nomenclature[Beta]{$\eta$}{Surtension (V)}
\nomenclature[Blambda]{$\lambda_0$}{Longueur de Debye (m)}
\nomenclature[Bnu]{$\nu_i$}{Coefficient stœchiométrique de l'espèce $i$}
\nomenclature[Bphi]{$\varphi$}{Module de Thiele microscopique}
%\nomenclature[Bphimacro]{$\varphi_\text{macro}$}{Module de Thiele macroscopique}
\nomenclature[Bphil]{$\varphi_l$}{Potentiel électrique dans le fluide (V)}
\nomenclature[Bphis]{$\varphi_s$}{Potentiel électrique dans le solide (V)}
\nomenclature[Bsigmas]{$\sigma_s$}{Conductivité électrique du solide ($\text{S m}^{-1}$)}
\nomenclature[Bsigmaq]{$\sigma_q$}{Densité de charge surfacique à l'interface ($\text{C m}^{-2}$)}
\nomenclature[Bsigmae]{$\sigma^e$}{Conductivité effective du milieu poreux ($\text{S m}^{-1}$)}

% ------------------------------------------------------------
% OPÉRATEURS ET MOYENNES
% ------------------------------------------------------------
\nomenclature[Cavg]{$\langle g \rangle$}{Moyenne superficielle de $g$}
\nomenclature[Cavgl]{$\langle g \rangle^l$}{Moyenne intrinsèque de $g$ dans le fluide}
\nomenclature[Ctilde]{$\tilde{g}$}{Fluctuation de $g$ à l'échelle du pore}